\documentclass[11pt]{article}

\usepackage[margin=1in]{geometry}
\usepackage[T1]{fontenc}
\usepackage[utf8]{inputenc}
\usepackage{lmodern}
\usepackage{setspace}
\usepackage{titlesec}
\usepackage{hyperref}

\setstretch{1.1}
\setlength{\parindent}{0pt}
\setlength{\parskip}{0.8em}

\titleformat{\section}{\large\bfseries}{\thesection}{1em}{}

\begin{document}

\begin{center}
    {\LARGE \textbf{UMass Lowell Summer 2026 KCS Research Proposal}} 
    
    \vspace{0.5em}
    
    {\large Jegannarayanan Sivakasiulaganathan, Student ID: 02197006}
    
    \vspace{0.25em}
    {\large Advisor: Dr. Paul Downen}
    
    \vspace{0.25em}

    {\large Graduate Student, Computer Science}
    
    \vspace{0.25em}
    
    {\large Expected Graduation: May 2027}
\end{center}

\vspace{1.5em}

\section*{Summary}
This Research Grant Proposal will focus on the implementation and proofs of
novel Dependent Type programming concepts.
Dependent types are types which depend on the values of their underlying terms.

I am interested in working with these types because I am interested in programming languages and type theory. 
I am interested in (and work on) compilers in my own time, and I am also currently taking a personal research course with Dr. Downen
supervising.


\section*{Planned Work}
The planned work is first to port new type structures that Dr. Downen and Shriya Thakur are working on
to Rocq. The work involved here is implementation - implementing the proposed types
in the proof solver Rocq. This is a good estimate of work for the 1st month.

Next, what is planned is to prove certain properties about these new types which make them desirable to be used.
Some of these are: soundness - essentially that the runtime guarantees that the type system proposes actually happen

As well as adequacy - ensuring all semantically correct operations succeed.


\section*{Expected Results}




\end{document}
